\chapter*{如何使用本手册}
\addcontentsline{toc}{chapter}{如何使用本手册}
\markboth{如何使用本手册}{如何使用本手册}

\section*{按你的处境查}

本手册按保研流程的时间顺序编排。如果你时间紧，可以按下面的方式直接跳到对应部分。

\begin{longtable}{p{4.6cm}p{4.0cm}p{7.6cm}}
\toprule
\textbf{你现在的情况} & \textbf{建议先看} & \textbf{原因} \\ \midrule
\endfirsthead
\toprule
\textbf{你现在的情况} & \textbf{建议先看} & \textbf{原因} \\ \midrule
\endhead
大一、大二，刚听说保研 & 第 1、2、4 章 & 先搞清制度、资格与排名怎么形成，再谈提升 \\
大三上，想知道自己够不够 & 第 2、3、4 章 & 资格与排名口径决定你能投哪里 \\
大三下，开始准备材料 & 第 5、6 章 & 材料与联系导师在这一阶段启动 \\
暑假，正在参加夏令营 & 第 7、10 章 & 夏令营规则与考核形式 \\
八月末，夏令营结果不理想 & 第 8、12 章 & 预推免是第二轮机会，同时要准备补救方案 \\
九月，推免系统即将开放 & 第 9、11 章 & 系统规则与确认时限最容易出错 \\
已拿到待录取，想确认后续 & 第 11、15 章 & 确认的不可逆性与入学前复审 \\
走专项计划（支教团、退役等） & 第 13 章 & 专项通道有独立的文件依据 \\
双非、边缘生、跨专业等 & 第 16 章 & 针对不同背景的专门问题 \\
想核实某个说法真假 & 第 19 章 & 常见误区逐条核对 \\ \bottomrule
\end{longtable}

\section*{每条问答怎么读}

每条问答以问题为题，答案在相关结论处标注数字引用。文末《参考文献》按同一数字编号列出发布机构、标题、年份和原始网址；点击正文编号可跳转查阅。政策和院校规则应以所引文件及当年正式通知为准。

\section*{关于“学校差异极大”的提醒}

有几个问题在本手册中反复出现，原因是它们在\textbf{不同学校、不同学院之间的答案差别太大}，任何统一说法都会误导人：

\begin{itemize}
\item \textbf{成绩排名的计算口径}（加权平均分、学分绩点、算术平均；含不含选修；重修如何计）——教育部不作统一规定，各校各院自定。
\item \textbf{英语门槛}——从“不设要求”到“六级 600 分”都有实例，且同一所学校的不同学院也可能不同。
\item \textbf{“优秀营员”的效力}——从“可直接拟录取”到“与录取完全无关”，跨度极大。
\item \textbf{待录取通知的确认时限}——短到几十分钟、长到数十小时，各校自定。
\item \textbf{加分项与竞赛认定目录}——各校的加分细则与认定清单不同，跨校套用经验容易出错。
\end{itemize}

碰到这五类问题，\textbf{本手册能给你的是“别人是怎么规定的”以及“你该去哪里查自己的规定”}，而不是一个可以照搬的结论。

\section*{书末的参考工具}

\begin{itemize}
\item \textbf{参考文献}：正文引用到的全部来源，按编号排列，含机构、类型、标题、年份、网址与被引次数。
\item \textbf{主题索引、阶段索引、院校与机构索引}：按主题词、保研阶段、院校机构三种方式检索全书。
\end{itemize}
